\documentclass[11pt]{report}

\usepackage{lmodern}
\usepackage[utf8]{inputenc}
\usepackage[T1]{fontenc}
\usepackage[french]{babel}

% package qui creer des liens dynamiques dans un fichier pdf a partir de tout element tagge 
%
\usepackage[linkcolor=blue,urlcolor=black,colorlinks=false]{hyperref}

% package qui permettent d'utiliser tout un tas de symboles mathematiques
\usepackage{amsmath} 
\usepackage{amssymb}


% package qui permettent d'inclure des graphiques
\usepackage{graphicx}
\usepackage{float}

% package qui permettent de personaliser les puces des listes
\usepackage{enumerate}

% package qui permettent d'introduire de la couleur
\usepackage{color}

\usepackage{xspace}

\usepackage{titlesec}

\titleformat{\chapter}[display]
  {\Huge\bfseries}
  {}
  {0pt}
  {Exercice \thechapter : \ }

\titleformat{name=\chapter,numberless}[display]
  {\Huge\bfseries}
  {}
  {0pt}
  {}


% sur une page Fige les parametres de hauteur et largeur du texte, de hauteur verticale au haut de page, de longueur horizontale au bord gauche de page de No pair, impair 
\setlength{\textheight}{24cm}
\setlength{\textwidth}{16cm}

\setlength{\topmargin}{-2cm}

\setlength{\evensidemargin}{0cm}
\setlength{\oddsidemargin}{0cm}

% Redefinit l'espace entre les lignes d'un paragraphe. Valur par defaut 1
\renewcommand{\baselinestretch}{1.2}

% Parametre l'espace apres un paragraphe
\parskip = 0mm
% Parametre l'espace d'indentation avant un debut de paragraphe
\parindent 5mm

% Positionne par rapport a la marge le premier niveau d'item d'une liste
\leftmargini 1.3em

% Ensemble et operateur avec un barre precedent le symbole
\newcommand{\E}{\mathbb{E}}     
\renewcommand{\P}{\mathbb{P}}     
\newcommand{\N}{\mathbb{N}}     
\newcommand{\Z}{\mathbb{Z}}
\newcommand{\R}{\mathbb{R}} 
\newcommand{\V}{\mathbb{V}} 
\newcommand{\T}{\mathbb{T}} 

% Lettres calligaphiques
\newcommand{\cA}{\mathcal{A}}
\newcommand{\cU}{\mathcal{U}}
\newcommand{\cN}{\mathcal{N}}

\usepackage{ntheorem}
% Definition des envrionement de type enonces numerotes simplement
\newtheorem{theo}{Théorème}
\newtheorem{rema}{Remarque}
\newtheorem{exem}{Exemple}
\newtheorem{defi}{Définition}
\newtheorem{exer}{Exercice}
\newtheorem{prop}{Proposition}
\newtheorem{lem}{Lemma}

% Definition des memes environnements de type enonces numerotes mais avec l'ajout prealable de la section 
\newtheorem{adef}{Définition}[section]
\newtheorem{apro}{Proposition}[section]
\newtheorem{athe}{Théorème}[section]
\newtheorem{aexe}{Exercice}[section]
\newtheorem{aexem}{Exemple}[section]
\newtheorem{arem}{Commentaire}[section]
\newtheorem{alem}{Lemma}[section]
%%%%%%%%%%%%%%%%%%%%%%%%%%%%%%%%%%%%%%%%%%%%%%%%%%%%%%%%%%%%%%%%%%%%%%%%%%%%%%%%%%%%%%%%%%%%%%%%%%%

\begin{document}

\chapter{Comparaison des algorithmes}

\section{Question 1 : }

The Schrödinger equation can describe the state of any quantum system: 
\begin{equation}
    t>0, \quad i\hbar\frac{\partial \Psi}{\partial t}(x,t) = \bigg(-\frac{\hbar^{2}}{2m}\frac{\partial^{2}}{\partial x^{2}}+V(x)\bigg)\Psi(x,t)=E\Psi(x,t)
\end{equation}

By using the separation of variables method we obtain the following expression for the wave function:
%\begin{figure}[H]
%\centering
%\includegraphics[scale=.4]{exo1/img/temp.png}
%\end{figure}

$$
\forall x, \forall t\geq 0: \Psi(x,t)=\phi(x) \nu(t) \Rightarrow \frac{\partial \Psi}{\partial t}(x, t)= \phi(x) \nu'(t)~ and~ \frac{\partial^2\Psi}{\partial x^2}(x, t)=\phi''(x)\nu(t)
$$

By substituting in the Schrödinger equation, we find that: 

$$
\underbrace{-\frac{\hbar^2}{2m} \frac{\phi''(x)}{\phi(x)}+V(x)}_{X}=\underbrace{ i \hbar \frac{\nu'(t) }{\nu(t)}}_{T}
$$

$X$ only depends on the space variable $x$, and $T$ only depends on the time variable $t$. And since $X=T$, we have $X=T=E$, with $E$ a constant that is homogeneous to an energy (same unit as $V(x))$.\\
By solving $i \hbar \frac{\nu'(t) }{\nu(t)}=E$, we have that $\nu(t)=A~ e^{-i\frac{E}{h}t}$, with $A$ a constant that is independent from $t$, so we can integrate it to $\phi(x)$.

\end{document}